\documentclass[11pt]{article}
\usepackage[margin=1.1in]{geometry}
\usepackage{amsmath, amssymb, amsthm}
\usepackage{booktabs}
\usepackage{xcolor}
\definecolor{linknavy}{RGB}{25, 55, 125}
\usepackage[colorlinks=true, linkcolor=linknavy, citecolor=linknavy, urlcolor=linknavy]{hyperref}
\usepackage{orcidlink}
\usepackage{fancyhdr}

\newtheorem{theorem}{Theorem}[section]
\newtheorem{proposition}[theorem]{Proposition}
\newtheorem{lemma}[theorem]{Lemma}
\newtheorem{conjecture}[theorem]{Conjecture}
\theoremstyle{remark}
\newtheorem{remark}[theorem]{Remark}

\newcommand{\C}{\mathbb{C}}
\newcommand{\Q}{\mathbb{Q}}
\newcommand{\Gm}{\mathbb{G}_m}

\title{Planar Keller Maps after the Counterexample:\\
Machine-Certified Newton-Polygon Exclusions,\\
the Staircase Transport, and the Recalibrated Frontier}
\author{Felipe Santiba\~nez-Leal\,\orcidlink{0000-0002-0150-3246}\\
\small CAOS Research program, \texttt{github.com/fsantibanezleal/CAOS\_RESEARCH}}
\date{2026-07-30\\[2pt]\normalsize Version 0.26}

\pagestyle{fancy}
\fancyhf{}
\fancyhead[L]{\small CAOS Research Preprint}
\fancyhead[R]{\small Planar Keller Maps \textperiodcentered\ v0.26}
\fancyfoot[C]{\small \thepage}
\renewcommand{\headrulewidth}{0.4pt}
\fancypagestyle{plain}{\fancyhf{}\fancyfoot[C]{\small \thepage}\renewcommand{\headrulewidth}{0pt}}

\begin{document}
\maketitle
\begin{center}\small
\fbox{\parbox{0.92\textwidth}{\centering
\textbf{PREPRINT} \textperiodcentered\ not peer reviewed \textperiodcentered\ CAOS Research programme\\
Paper B of a three-paper record \textperiodcentered\ version 0.26 \textperiodcentered\ 2026-07-30\\
DOI \href{https://doi.org/10.5281/zenodo.21503367}{10.5281/zenodo.21503367} (concept DOI, resolves to the latest version)\\
Licensed CC BY 4.0 \textperiodcentered\ sources and machine records: \texttt{github.com/fsantibanezleal/CAOS\_RESEARCH}
}}
\end{center}


\begin{abstract}
The July 2026 counterexample settled the Jacobian conjecture in dimensions $N \ge 3$ and left
$JC(2)$ as the open case. This manuscript records an exact-arithmetic program on the planar
case. The Keller condition $J(P, Q) = 1$ is linear in $Q$ for fixed $P$; the resulting machine
(eliminate, parametrize the consistency variety, complete, invert) closes the columns
$\min\deg \le 2$ uniformly (one closed inverse for the whole quasi-triangular class), forces
perfect-power tops at bidegrees $(2, n)$ and $(3, n)$, and inverts every consistent stratum it
has met. Four unconditional exclusion theorems follow from a weight-class calculus: a two-term
$P = x + a x^u y^v$ is never a Keller component at any partner degree; neither is any
lower-weight perturbation of it; no Keller component carries a Newton top edge anchored at the
linear vertex $x$ with interior direction; and if $x$ is a vertex of the Newton polygon then
$P = x + f(y)$ exactly (the vertex dichotomy). The residual open core is the transport of the
Keller constant along the staircase of Newton vertices swallowing $x$: we show the completion
window is block-triangular under the edge grading, exhibit the classwise transport chain as an
executable elimination whose obstruction always localizes at the constant's class, and prove the corresponding
fifth exclusion: for primitive tops the certificate-tower induction closes it
UNCONDITIONALLY at all partner degrees (cleared base certificates with monomial pairings
such as $-13860\,a^7$, plus a resonance-reduction lemma showing window growth can never
repair the constant's row). A source-audited recalibration places the program
honestly: our composite-gcd certificates replicate literature-covered territory (the verified
floor is $\gcd \le 8$, primes, $2p$, then Heitmann's $B \ge 16$ and the
Guccione--Guccione--Valqui dichotomy $B = 16$ or $B > 20$), and the live frontier is the
$B = 16$ normal form together with the single surviving degree pair $(72, 108)$ below $125$.
For that open pair, an exact primary-source instantiation gives two original-pair rejection
gates: seven possible northeastern vertices for a nonzero inner polynomial, and
$\deg_x\operatorname{Res}_y(P,Q)=21$ with an $84+24$ major/minor root partition; neither
gate alone excludes the pair.
A constructible determinantal-strata calculation now closes the first actual
three-coefficient GGHV slice exactly. On the residual curve, two normalized
maximal minors reduce, with \(z=u^9\), to
\(21-96z-1024z^2\) and \((8z+1)^{14}\); a persisted integer Bezout identity
has constant \(17^{14}\). After adding \(\varepsilon_{(0,7)}\), the only
two-chart residual is a degree-12 fiber over \(z=-1/8\). A third exact minor
has coprime degree-13 fiber determinant, certified by a second integer
Bezout identity. A connected mod-9 row/column grading extends to every
nonconstant direction in the persisted lower family. This is a strict slice
result and does not exclude the full \((72,108)\) family.
The full-family audit then removes the structural constant column, corrects the
complete row count from 289 to 302, and compresses one selected augmented
minor to a 36-column block depending on 24 of the 51 parameters. Two
connectivity-generating triples factor exactly; one graph-active direction
cancels identically, while the other determinant becomes weighted homogeneous
with weights \((7,3,9)\) after shifting the forced coefficient. On the
\(d=0\) boundary, the tautological kernel \([P,P]=0\) can be quotiented
exactly. One quotient determinant factors into \(b^{32}\), one five-term
weighted factor, and five binomials; a second is a nonzero scalar times
\(a^{107}\). Together with the exact origin rank gap, these charts close the
complete quotient boundary. On \(d\ne0\), a first alternative chart reduces
each selected component \(G,L,Q\) to a finite intersection. A third exact
determinant factors as \(X^{30}LQR\), for \(X=A^3\), and exact component
ideals close \(G\). A row basis selected independently on the finite \(L\)
and \(Q\) residuals is in fact shared. Its 23-term determinant descends as
\(A^{87}R(A^3,B)\); exact unit-ideal certificates close both residuals.
Together with the quotient boundary this closes the complete
three-parameter \(T_B\) restriction. This is still a strict slice result,
not closure of the 24-parameter core, the full family, \((72,108)\), or
\(JC(2)\).
Restoring direction \((2,9)\) gives the first higher-dimensional recursive
atlas. The shared determinant is
\(A^{87}(R(X,B)+YS(X,B))\), \(X=A^3\), \(Y=A^2C\), with
\(\gcd(R,S)=1\). A second exact minor covers a dense open of this rational
graph; its residual factors into three plane curves \(F_3F_6F_7=0\).
On the linear-in-\(X\) component
\(F_3=(5B+4)^3+16X\), a third exact minor reduces the principal-open residual
to \(Q_9(B)Q_{15}(B)=0\): 24 normalized values, or 72 lifted algebraic
points. On the quadratic-in-\(X\) component \(F_6\), the same exact section
has a nonzero linear quotient class \(U(B)X+V(B)\). Its degree-74
function-field norm has five irreducible factors of degrees \(2,3,6,18,30\).
Exact same-point arithmetic places the first three on \(X=0\) or \(R=S=0\);
only the degree-18 and degree-30 factors survive on \(AS\ne0\). Thus \(F_6\)
also becomes finite: 48 normalized values, or 144 lifted algebraic points.
The curve \(F_7\), the finite base locus, and \(A=0\) remain.
A new multiplication formula for $x^m$-anchored edges shows the $B = 16$ leading form cannot
produce the Keller constant on its top edge for $m \ge 2$: that case, too, is a staircase
transport. All computations are exact and reproducible; every claim is labeled machine-verified
[MV], derived [D], or conjectural [C].
\end{abstract}

\section{Introduction}

Let $P, Q \in \C[x, y]$ with $J(P, Q) := P_x Q_y - P_y Q_x \in \C^\times$ (a \emph{Keller
pair}; each member is a \emph{Keller component}). The two-variable Jacobian conjecture $JC(2)$
asserts every Keller pair is an automorphism pair. Following the July 2026 counterexample in
dimension 3 \cite{Alpoge2026}, $JC(2)$ is the surviving case, and the classical planar theory
(Newton polygons, approximate roots, gcd coverage) is its context. This manuscript is the
planar half of a two-manuscript record: the foundational companion documents the
three-dimensional aftermath (validation, the seed family, the geometry of non-properness,
two-dimensional equivariant rigidity); this one documents the planar program. The split
preserves history: every statement cites its experiment record (EXP-$nnn$) in the repository,
where hypotheses are declared before runs and refuted attempts are kept.

Labels: \textbf{[MV]} machine-verified exact identity; \textbf{[D]} derived analytically and
certified on instances; \textbf{[C]} conjecture. Machine-verification is exact sympy
arithmetic over $\Q$ or $\Q(i)$; no floats anywhere.

Two structural inputs from the companion are used freely. First, every $\Gm$-equivariant
Keller map of $\C^2$ is an automorphism; on the opposite-sign/one-zero scope relevant to the
three-dimensional mechanism it is linear (EXP-010; independently classified for all signatures
by Shaska \cite{Shaska2026}). Second, at every weight ray the leading pair of a planar
Keller map is Jacobian-dependent (the tropical bridge, EXP-013), which is the entry point of
all edge analyses below.

\section{The planar machine}

The Keller condition is bilinear in the coefficient vectors of $P$ and $Q$, hence
\emph{linear in $Q$ given $P$}. In the affine gauge (identity linear parts) this yields a
machine per degree pair: eliminate $Q$ (the consistency ideal in $P$'s coefficients),
parametrize, complete linearly, and invert or fiber-test (EXP-017, EXP-019, EXP-020).

\begin{theorem}[{[MV]}; uniform closure at $\min\deg \le 2$ (EXP-021/022)]\label{thm:uniform}
Write $\ell = \alpha x + \beta y$, $\beta \ne 0$. In shear coordinates the Keller condition
for $P = x + f(\ell)$ (any polynomial $f$) is a linear first-order PDE with characteristic
invariant $P$; its complete polynomial solution space is $Q = \ell/\beta + H(P)$, $H$
arbitrary, and the pair inverts by one closed formula. Consequently every planar Keller map
with $\min(\deg P, \deg Q) \le 2$ is invertible, and the whole quasi-triangular class
$P \sim x + f(\ell)$ is invertible at every degree.
\end{theorem}

\begin{theorem}[{[MV]}; forced tops at $(2, n)$ and $(3, n)$ (EXP-019/020/022/039)]
\label{thm:forced}
The $(2,3)$ consistency ideal is the single equation $(4A_0A_2 - A_1^2)^2 = 0$ ($P$'s top
form is a perfect square); the $(2,4)$ elimination returns the same ideal. At $(3, 4)$, the
$6 \times 5$ linear map $Q_4 \mapsto J(P_3, Q_4)$ is rank-deficient exactly when $P_3$ is a
perfect cube: all $5 \times 5$ minors vanish on the cube parametrization, every minor lies in
the Hessian-covariant ideal, and completeness follows from $GL_2$-equivariance
($J(f \circ A, g \circ A) = \det(A)\, J(f, g) \circ A$, verified identically) plus full rank
at the two non-cube orbit representatives $x^2 y$ and $x y (x + y)$. On the cube stratum,
sampled completions are Keller and invert with explicit polynomial inverses ($4/4$ by lex
Gr\"obner extraction). At $(3,3)$ the consistency ideal forces the quasi-triangular alignment
at the radical level.
\end{theorem}

The degree columns close by classical coverage: $\gcd(2, n) \le 2$ and $\gcd(3, n) \in
\{1, 3\}$ are inside the verified $\gcd$ landscape of Section~\ref{sec:frontier}, so the
machine's contribution there is the \emph{constructive} loop (explicit consistency varieties
and inverses), stated as replication. The classical degree descent is an algorithm on this
machine: subtract $c P^k$ while the tops are power-proportional, finish with
Theorem~\ref{thm:uniform}; it explicitly inverted a deterministic library of composed Keller
maps, and rotate-descent (rotation of linear-base tops plus de Jonquieres subtraction,
EXP-034) linearizes that entire library in $0$ to $6$ steps. Its exact limit: mixed-corner
tops stay mixed under rotation; those are the classical hard shapes, and they are where the
exclusion theorems below take over.

\section{Four unconditional exclusions}

Throughout, $P$ has linear part $x$ after affine gauge. The weight-class calculus: under
$w = (v, 1 - u)$, the operator $J(P, \cdot)$ shifts weight by a constant, so the Keller
equation decouples into weight classes; the technique goes back to Dixmier and Joseph
\cite{Joseph1975} and is standard in the polygon literature \cite{GGV2017}; the statements
below appear to be new in the stated generality (Section~\ref{sec:positioning}).

\begin{theorem}[{[MV/D]}; the weight-class theorem (EXP-029)]\label{thm:t1}
For $u \ge 2$, $v \ge 1$, $a \ne 0$, the polynomial $P = x + a\,x^u y^v$ is not a Keller
component: no $Q$, of any degree, satisfies $J(P, Q) \in \C^\times$.
\end{theorem}

\begin{proof}[Proof shape]
$P$ is $(v, 1-u)$-homogeneous; the constant is reachable only from the class of $y$, spanned
by the ray $g_s = x^{ks} y^{ms+1}$ ($k = (u-1)/d$, $m = v/d$, $d = \gcd(u-1, v)$); there the
operator is banded, $J(P, g_s) = (ms+1)\,e_s + a B_s\,e_{s+d}$ with $B_s$ a positive integer,
and the chain recursion contradicts the last row of every truncation.
\end{proof}

\begin{theorem}[{[MV/D]}; lower-weight perturbations never rescue (EXP-030/031)]\label{thm:t2}
With $u, v, a$ as above and $R$ any polynomial whose monomials have degree $\ge 2$ and
$(v, 1-u)$-weight strictly below $v$: $P = x + a\,x^u y^v + R$ is not a Keller component.
\end{theorem}

\begin{theorem}[{[MV/D]}; every $x$-anchored edge falls (EXP-032)]\label{thm:t3}
Let $z = x^k y^m$, $k, m \ge 1$, and $E = x\,\varphi(z)$, $\deg\varphi = g \ge 1$. On the
$y$-class the edge operator is \emph{multiplication}:
$J(E, g_s) = [(ms + 1)\varphi + k z \varphi'] z^s$, and the top coefficient
$(mD + kg + 1)\varphi_g$ of the class ODE kills every polynomial candidate. Hence $E$ is
never the leading form of a Keller component, and $P = E + R$ (any strictly lower-weight
tail) is never a Keller component.
\end{theorem}

\begin{theorem}[{[MV/D]}; the vertex dichotomy (EXP-033)]\label{thm:t4}
If the monomial $x$ is a vertex of the Newton polygon $N(P)$ (convex hull of the support
with the origin) of a Keller component with linear part $x$, then $P = x + f(y)$.
\end{theorem}

\begin{lemma}[Annihilation (EXP-036)]\label{lem:annihilation}
For any polynomial $P$, $J(m, P^k) = -k\,L(P^{k-1} m)$ with $L = J(P, \cdot)$ (Leibniz). Any
covector annihilating the image of $L$ on a space containing the preimage therefore kills the
source; the certificate covectors used throughout are left-null on the completion matrix,
whose columns are exactly $L(m)$, so the tail steps of
Theorems~\ref{thm:t2}--\ref{thm:t4} are unconditional. The finite-window subtlety (spurious
pairings for out-of-window preimages) reproduces the recorded artifact values exactly.
\end{lemma}

The dichotomy's other side is real: $x + (x+y)^2$ and $x + (x+2y)^3$ are components in which
$x$ is \emph{swallowed} inside the polygon (equivalently: a pure power $x^d$, $d \ge 2$,
appears in the support). Every non-triangular component is of that kind. The residual open
core is therefore precise: components whose polygon swallows $x$ with a mixed staircase.

\section{The corner is diagonal; the staircase transport}\label{sec:staircase}

\begin{proposition}[{[MV]}; the diagonal corner (EXP-035)]
For $B = x^i y^j$, $i, j \ge 1$:
$J(B^p, x^a y^b) = p (ib - ja)\, x^{ip + a - 1} y^{jp + b - 1}$, with kernel exactly the
powers of $B$; no monomial gives $J(B^p, \cdot)$ a constant term. The mixed corner is clean;
the obstruction lives on the staircase of Newton vertices joining the swallowed linear
vertex to the corner.
\end{proposition}

\begin{theorem}[{[MV]}; block triangularity and the transport chain (EXP-037)]
\label{thm:transport}
Fix the $x$-edge grading $w = (v, 1-u)$ and let $s_0$ be the minimal $w$-class of $P$. In
the completion window (unknown $Q$ of bounded degree), every nonzero matrix entry sits at a
class offset $s - \sigma$ for a $P$-class $s$ ($\sigma = v + 1 - u$), and the offset
$s_0 - \sigma$ is minimal: ordering $Q$-classes by $w$-class, the system is block
triangular. Classwise elimination (solve the diagonal block, inject its kernel parameters,
emit its cokernel obstructions) is therefore well-defined; it reproduces the monolithic
consistency verdict on every tested sample, including a consistent positive control. On the
pure above-weight staircase family the diagonal blocks are the banded operators of
Theorem~\ref{thm:t1}, and on all $8$ grid samples the first inconsistent transport equation
sits at the \emph{constant's} class.
\end{theorem}

\begin{theorem}[{[MV]}; the fifth exclusion, window form (EXP-042)]\label{thm:t5}
Let $P = x + a\,x^u y^v + b\,x^d$ ($u \ge 2$, $v, d \ge 1$; the minimal shape swallowing
$x$). For every grid point $(u, v, d) \in \{2,3\} \times \{1,2\} \times \{2,3\}$ and window
$N = 7$, a cleared certificate covector with polynomial entries satisfies $c^T M = 0$
identically and pairs with the right-hand side to a \emph{monomial}:
\[
-13860\,a^7,\quad -32\,a^4,\quad -630\,a^4,\quad -2730\,a^7,\quad -1155\,a^3,\quad
-165\,a^4,\quad -4620\,a^7 b,\quad -81\,b^4,
\]
so the completion window is empty for \emph{all} parameter values with $a \ne 0$ (the
$b$-factors vanish only where Theorem~\ref{thm:t1} applies anyway). At $(u, v, d) =
(2, 1, 2)$ the pairing at window $N$ is $-c_N\,a^N$ with $c_N \in \{420, 1260, 13860,
180180, \dots\}$ for $N = 5..9$: a nonzero monomial in $a$ alone at every certified window.
The annihilation mechanism of Lemma~\ref{lem:annihilation} transfers verbatim (it is
$P$-agnostic), and in-window sources pair to zero exactly as the window criterion demands.
A three-parameter tail certificate ($-27720\,a^7 c_3^4$) extends the family.
\end{theorem}

\begin{lemma}[{[MV/D]}; the tower lemma (EXP-044/046)]\label{lem:tower}
Let $P = x + R$ ($R$ without constant or linear part), $M_N$ the completion-window
matrix at window $N \ge \deg P - 1$, and $T$ the total-degree top form of $P$. Suppose
every form in $\ker J(T, \cdot)$ at each degree is a scalar multiple of a power that
reduces, modulo $\C[P]$, to degree $\le N$ (this holds whenever $T$ is a primitive
monomial, and whenever $T = P - x$ or $T$ is the pure power $b x^d$, by expanding
$T^k = \lambda P^k + \text{lower}$). Then any left-null certificate of $M_N$ with
nonzero pairing extends to $M_{N+1}$ with the same (scaled) pairing: window
inconsistency propagates to every larger window. Proof ingredients, each
machine-verified: old columns vanish on new rows (degree bookkeeping); the new block is
$J(T, \cdot)$ on degree-$(N{+}1)$ forms with the predicted kernel; each resonance
$\kappa$ satisfies $\kappa - \lambda P^k \in$ window with $P^k \in \ker L$, so
$L(\kappa)$ is an in-window image which every certificate annihilates; the extension
was realized explicitly across an active resonance (window $8 \to 9$, ratio $2$,
pairing preserved). The measured window law ($-c_N a^N$ pairings with the odd-prime
$2N - 3$ ratios) is this tower seen through the clearing normalization.
\end{lemma}

\begin{theorem}[{[MV/D]}; all-degree fifth exclusion (EXP-042/044/046)]
\label{thm:t5all}
For $u \ge 2$, $v \ge 1$ with $x^u y^v$ primitive ($\gcd(u, v) = 1$) or $d \ge u + v$,
and any $a \ne 0$, $b$: $P = x + a\,x^u y^v + b\,x^d$ is not a Keller component at ANY
partner degree. (Base: the cleared certificates of Theorem~\ref{thm:t5}; induction:
Lemma~\ref{lem:tower}.) For proper-power tops (e.g. $u = v = 2$) the odd resonances do
not reduce mod $\C[P]$; the certificates kill them anyway in every tested instance
(e.g. $(xy)^5$ at window $9$), so the statement there is [D] pending the derivation of
that kill.
\end{theorem}

\begin{theorem}[{[MV/D]}; the universal tower (EXP-049)]\label{thm:t6}
Let $P = x + R$ ($R$ without constant or linear part) whose total-degree top form
$T = P_D$ is NOT a proper power. If the completion window is inconsistent with a
certificate at any single $N \ge \deg P - 1$, then $P$ is not a Keller component at ANY
partner degree. (The tower argument with $T$ the full top form: $\ker J(T, \cdot)$ on
degree-$M$ forms is zero unless $\deg T \mid M$ and then the span of $T^{M/\deg T}$;
and $T = P - \text{lower}$ by definition, so $T^k - P^k$ has degree $\le k \deg P - 1$:
every resonance reduces in-window modulo $\ker L$.) One cleared window solve, seconds
of compute, now yields an all-degree exclusion per shape: the machine harvest includes
multi-edge staircases such as $x + a x^2 y + b x^2 + c x^3 y^2$ (pairing $-4620\,c^4$).
For proper-power tops the measured mechanism is the HALF-PLANE certificate (EXP-048): a
pairing-nonzero certificate supported in the $y$-heavy half-plane, disjoint from every
$x$-heavy resonance source; its general construction is the program's remaining
structural gap, and it is exactly what the $B = 16$ frontier shapes (proper-power tops)
require. First-contact certificates now exist at the open pair's degrees themselves:
sampled degree-72 shapes with the GGHV corner $(16, 56)$ have empty filtered windows at
$N = 108$ in about one second each (EXP-050).
\end{theorem}

\begin{theorem}[{[MV/D]}; the half-plane tower (EXP-051)]\label{thm:halfplane}
Let $P = x + R$ whose total-degree top form attains $\min\{p - q\}$ over the support
(the swallowed-staircase geometry), and let $H$ be the $y$-heavy half-plane of rows
$\{i \le j\}$. In the $H$-restricted window system every tower case resolves: $T$-power
resonances reduce by the universal identity; other kernel resonances have $x$-heavy
sources with empty $H$-part; columns whose $T$-output leaves $H$ vanish from the
subsystem. Hence ONE $H$-certificate at one window excludes $P$ at ALL partner degrees,
proper-power tops included. Frontier payoff: the degree-32 $B = 16$-flavor sample and
the degree-72 corner-$(16, 56)$ sample carry $H$-certificates (pairings $-256$, $-32$)
and are excluded at all partner degrees.
\end{theorem}

Two instrument-level results complete the toolkit. The \emph{matched-pair law} (EXP-038):
for cornered pairs $P = \alpha B^p + P_{\mathrm{sub}}$, $Q = \beta B^q + Q_{\mathrm{sub}}$,
the second class of $J(P, Q)$ is exactly
$\alpha p (ib - ja) q_{a,b} - \beta q (id - jc) p_{c,d}$ over outputs matched by
$(a, b) = (c, d) + (q - p)(i, j)$; imposing the pair shape does not change the obstruction
depth on any tested sample, but it halves the window unknowns: an efficiency tool, not a
second obstruction source. And the \emph{$x^m$-anchored edge operator} (EXP-043): for
$E = x^m \varphi(z)$, $z = x^k y^l$,
\[
J\left(x^m \varphi(z),\; x^{\alpha} y^{\beta}\right)
= x^{m + \alpha - 1}\, y^{\beta - 1}\,
\bigl[\beta m\, \varphi(z) + (\beta k - \alpha l)\, z\varphi'(z)\bigr],
\]
verified identically in symbolic edge coefficients; Theorem~\ref{thm:t3}'s operator is the
case $m = 1$.

\section{The frontier, recalibrated from primary sources}\label{sec:frontier}

An audited literature pass (2026-07-22; dossier in the repository with per-claim sources and
explicit UNVERIFIED flags) fixes the honest coordinates of the program.

\emph{Verified gcd coverage.} For a counterexample pair, $B := \gcd(\deg P, \deg Q)$
satisfies: $B \ge 9$ is forced by the classical interval results ($B = 1$: Magnus
\cite{Magnus1955}; $B \le 2$: Nakai--Baba; $B \le 8$ or prime: Appelgate--Onishi
\cite{AppelgateOnishi1985} and Nagata \cite{Nagata1989}); $B \ne p$ \cite{Abhyankar2008};
$B \ge 16$ (Heitmann \cite{Heitmann1990}); $B \ne 2p$ and $B = 16$ or $B > 20$
\cite{GGV2017}. Zoladek's claimed $B > 33$ has a documented gap \cite{GGV2017} and is not
used. Consequently our machine certificates at $\gcd = 2, 9, 12, 18$ (EXP-024..028), stated
in earlier revisions as frontier contact, are \emph{replications} of literature-covered
territory; we keep them as independent verifications and relabel them accordingly.

\emph{Verified degree floor.} Moh's classical analysis discards $\max\deg \le 100$
\cite{Moh1983}, with complete published detail only for the largest case ($64$); the
four exceptional cases correspond to six configurations in the modern corner terminology,
which the algorithmic audit of Guccione--Guccione--Horruitiner--Valqui enumerates and
discards \cite{GGHV2017, GGHV2022}. The current floor is: any counterexample has
$\max\deg \ge 125$, or $(\deg P, \deg Q) \in \{(72, 108), (108, 72)\}$ \cite{GGHV2022};
that single pair was left open explicitly for compute reasons.

\emph{The live targets.} (i) The $B = 16$ case comes with a forced direction set and
leading forms $l_{4,-1}(P) = \lambda_P R_0^m$, $R_0 = x (x y^4 - \lambda_0)^3$, and a
reduced normal form \cite{GGV2017}. We verify [MV] that $R_0$ is $(4,-1)$-homogeneous with
collinear $x$-anchored support, and, by the $x^m$-anchored operator above, that for
$m \ge 2$ \emph{no monomial reaches the Keller constant from the $R_0^m$ edge}: the
$B = 16$ problem is a staircase transport problem, the exact object of
Section~\ref{sec:staircase}. The pure $m = 1$ shape has $x$ as a Newton vertex and is
excluded outright by Theorem~\ref{thm:t4} (window replication empty); swallowed variants
land in Theorem~\ref{thm:t5}'s territory (sampled windows empty). (ii) The similarity filter
(partner supported in the scaled polygon $\tfrac{\deg Q}{\deg P} N^0(P)$; sound by the
similarity theorem since the scaled polygons are nested in the degree) cuts the
$(48, 64)$ full-window system from $2142$ to $111$ unknowns, and the validation runs are
DONE: three degree-48 samples of the $F_1$ flavor (the pure shape $x + R_0(1)^3$ and two
swallowed variants) have EMPTY filtered windows at $N = 64$ in $2.4$ to $3.6$ seconds
each, excluding all partners of degree $\le 64$ per sample (the classical case-64
degrees: a sampled-level replication of Moh/Heitmann demonstrating the pipeline's cost
at target scale). (iii) For the $(72, 108)$ system the histograms give $5992$ unknowns
in $535$ classes with largest diagonal block $22$ before filtering: with the filter and
the transport, the open territory (max degree $> 150$, and $(72, 108)$ itself) is
affordable at seconds-to-minutes per certificate.

\subsection*{The $(72, 108)$ attack: the reduced system and the perturbative covector}

The lone surviving degree pair below $125$ reduces (Guccione--Guccione--Horruitiner--
Valqui \cite{GGHV2022}, Prop.~4.3; transcription verified against the LaTeX sources)
to the existence of a pair with $[P, Q] = x^2$ supported in SMALL polygons:
\[
\begin{aligned}
N(P)&:\ (0,0),(1,0),(8,14),(8,16),(0,8) &&\text{(61 lattice points)},\\
N(Q)&:\ (0,0),(2,1),(12,21),(12,24),(0,12) &&\text{(125 lattice points)}.
\end{aligned}
\]
The original paper's unsolved system was never printed. Machine results to date (EXP-050 through
EXP-055): sampled shapes at the original degrees and on the reduced polygons are all
partner-free (seconds per certificate); on the stratum with the forced top edge
$y^8 (xy - t)^8$ symbolic, every cleared certificate pairs to the CONSTANT $576$
(all $t$, identical across sampled lower strata; five-row support); rigid universality
of that covector was REFUTED by the machine (51 entering lower combinations), which
pointed to the perturbative construction $\Lambda(\varepsilon)$; its base is sound
(the right-hand side is the constant $x^2$ vector; $\Lambda_0$ globally left-null,
symbolic $t$) and FIRST-ORDER UNIVERSALITY is proved: all 26 entering lower monomials
admit correctors with the $576$ pairing pinned and localized supports. Order-1
termination fails (measured), so the corrector ladder continues at order 2; each rung
is a batched linear solve, and termination at any finite order yields the explicit
universal covector, closing the reduced stratum for all parameters. The remaining
assembly after that: the reduction's other forcing branches, then the floor rises from
$108$ to $125$.

\subsection*{The closure of the $(72, 108)$ case (validated 2026-07-22)}

The endgame completed: all three branches of the GGHV reduction are certificate-covered
on their forced families. Cases a) and b) (the sub-polygon systems) close with monomial
pairings on two normalization charts ($200\,a_0^3$, $600\,a_0^3$, $560\,a_0^4$;
$200\,a_1^4$), whose vanishing loci sit exactly where those polygons' own vertices
would degenerate, excluded by their forcing; the a/b partner side is covered by the
bigger-polygon emptiness. Case c) closes under the torus gauge (the forced-edge
parameter normalizes to $t = 1$; verified concretely by transporting an orbit sample
with identical verdict): the $\beta$-symbolic certificates give pairing
$23592960\,\beta$ with the degenerate $\beta = 0$ corner covered by the stratum
certificates, and the interior coefficients certify axis-symbolically with constant
pairings ($368640$ at every tested point), on top of dense mixed-direction sampling.
The orientation swap is absorbed by $Q \to -Q$. AMENDED after our commissioned adversarial audit
(same day): since the GGHV reduction forces NO interior coefficients, the closure must
hold for ALL interior values, and axis-symbolic-plus-sampled coverage does not yet
establish that (our own EXP-054 refutation is the in-house proof that slice
extrapolation fails). The certified statement is therefore: \textbf{the three reduced
branches are certificate-covered on their forced-edge families, with interior
generality axis-symbolic and densely sampled: strong evidence toward, not yet a proof
of, the $(72, 108)$ closure}; the floor-raise to $125$ becomes assertable only after
the simultaneous-symbolic interior certificate (the ranked hardening queue is in
\texttt{program/jacobian-conjecture/72108-closure-statement.md}).

\subsection*{The truncation ladder for the simultaneous certificate (2026-07-24)}

The certified statement above reduces the floor-raise to one object: a covector
$\Lambda(\varepsilon)$, polynomial in the interior coefficients $\varepsilon$,
that closes the reduced system for \emph{all} interior values at once. We report
the machine campaign that decides this object degree by degree; every step is
exact, and every modular reduction of a rational uses $\mathrm{num}\cdot
\mathrm{den}^{-1} \bmod p$ (a bug that truncated instead, see below, was caught by
a declared regression gate).

\emph{Solvability is automatic (the dual annihilation identity, EXP-068).} The
reduced system $M_0$ has corank one; its cokernel covector $\Lambda_0$ satisfies
$\Lambda_0 M_0 = 0$ and, more, it annihilates every Jacobian-bracket image that
arises in the corrector recursion. Hence the consistency (solvability) conditions
of the ladder hold at every order, and the only obstruction to a finite covector
is the top-order \emph{vanishing} condition. This is the dual form of the
annihilation lemma $J(m, P^k) = -k\,L(P^{k-1} m)$ of Section~\ref{sec:staircase}.

\emph{The tower, decided through degree two.} A polynomial covector of
$\varepsilon$-degree $d$ exists iff a finite linear system over the interior
coefficients is feasible. We find [MV]: \textbf{no covector of degree $1$}
(EXP-067; eight ``diagonal'' operators obstruct, the full system infeasible modulo
two independent primes, hence over $\mathbb{Q}$); and \textbf{no covector of
degree $2$} (EXP-072; the support $\{(0,1),(1,0),(3,5)\}$ obstructs, again
conclusive modulo two primes). At degree $3$ no obstruction is found through the
pair and all triple supports (EXP-071, EXP-073); the quadruple tier and a
constructive solve are in progress, subject to the one-sidedness noted below. The
obstruction \emph{persists} rather than migrating: the degree-two obstructing
support $\{(0,1),(1,0),(3,5)\}$ contains two of the eight operators that already
obstruct at degree one, and what grows is the support size needed to expose it
(single operators at degree one, a triple at degree two). An earlier revision of
this manuscript reported instead that the obstruction \emph{moves} to an unrelated
mixed pair; that reading rested on the retracted computation described next and is
corrected here.

\emph{An honesty note (EXP-070, retracted).} A first ``degree $2$ closed'' result
was withdrawn: its modular assembly reduced rationals by integer truncation, and a
regression gate declared before the next experiment refused to reproduce a known
exact decision, exposing the bug. Degree two was then re-decided soundly at a
different support. We keep the retraction in the record; the declared-gate
discipline is part of the method.

\emph{One structural reduction, and one claim withdrawn.} The rational-covector
relaxation adds nothing (EXP-074): since $\Lambda_0 M_0 = 0$ the certificate
condition is homogeneous, so any rational covector clears its denominator to a
polynomial one. A previous revision of this manuscript (v0.09, v0.10) asserted in
addition that the tower has a \emph{computable finite ceiling}, namely the
dimension of the Krylov closure of $\Lambda_0$, and concluded that the campaign is
a finite decision procedure. \textbf{That assertion is withdrawn here} (EXP-078):
its derivation presupposes that the pinned corrector ladder terminates, and
EXP-064 measured that it does \emph{not} (the descending chain stabilizes nonzero
at dimension $39$). Since the pinning carries a large gauge freedom, neither the
ceiling nor its negation follows; we currently have \emph{no} a priori degree
bound, effective or otherwise, and the chain of solvable degrees is controlled
only by Noetherianity, which is not effective.

What survives is weaker but exact. The set of degrees at which a covector exists
is \emph{upward closed} (a degree-$d$ covector is a degree-$(d+1)$ covector with
vanishing top coefficients), so either none exists at any degree, or there is a
minimal $d_0$; with degrees $1$ and $2$ closed, $d_0 \ge 3$. And the
support-restricted sweeps are \emph{one-sided}: they test necessary conditions, so
an infeasible support decides a degree negatively, while all supports passing does
\emph{not} establish existence. Settling degree three positively therefore
requires constructing a covector, not sweeping. Both readings of the outcome remain
open and we state them plainly: a covector would be an exclusion proof, and a proof
that none exists at any degree would be the first step of a counterexample.

\emph{A lead toward a degree-uniform statement.} Written as a connection, the
annihilation lemma makes the covector obstruction a first cohomology class, and
termination of the corrector ladder becomes the regularity type of that connection
(regular singular: the ladder terminates and a structural degree bound follows;
irregular: no such bound, matching the measured non-termination of the pinned
ladder). Whether this decides $(72,108)$ in one stroke is the current deep
question; it is stated as a program item, not a result.

\subsection*{Original-pair source gates for the open chain (EXP-095/096; validated 2026-07-25)}

The reduced Proposition~4.3 system is a Laurent pair with bracket $x^2$, so
theorems requiring a polynomial Keller pair with bracket one cannot be applied
to it directly. Returning instead to the hypothetical original pair gives an
exact applicability bridge [MV]. Its degree-$72$ component has
\[
  2A_0=(16,56),\qquad 2A'_0=(2,0),\qquad
  2A_1=(11/2,14),
\]
which is precisely the first retained $D=72$ branch in the Newton-resolution
list of Makar-Limanov--Trakhtenberg \cite{MakarLimanovTrakhtenberg2024}.
Thus that independent classification retains the GGHV branch; it does not
exclude it.

Two sharper necessary conditions do result. First, instantiating
Lee--Li's inner-polynomial region \cite{LeeLi2024} with
\[
 (a,b,m,n)=(2,3,16,56),\qquad \mathfrak m=\frac{139}{39},
\]
gives [MV] exactly seven possible northeastern vertices whenever the inner or
innermost polynomial is nonzero:
\[
 (1,3),(2,7),(3,10),(4,14),(5,17),(6,21),(7,24).
\]
The zero-inner-polynomial alternative remains. Nor does the stronger diagonal
corollary apply, since
\[
 \frac ba=\frac32\not>\frac{n-m}{a}-1=19.
\]

Second, the approximate-root formulas of
Guccione--Guccione--Horruitiner--Valqui \cite{GGHVApproxRoots} become an exact
intersection invariant [MV]. For the complete chain
\[
 (m,n)=(3,2),\qquad A_1=(11/4,7),\qquad k=1,\qquad l=4,
\]
the forced fourth-power edge supplies four final major classes. Each has
$|D_\tau^P|=3\cdot7=21$ and $\lambda_\tau^Q=1/4$, hence
\[
 I(P,Q)=\deg_x\operatorname{Res}_y(P,Q)
       =4\cdot21\cdot\frac14=21.
\]
The degree-$108$ component therefore has $84$ major and $24$ minor roots.
These are original-pair rejection gates independent of sampling the $51$
interior coefficients of the reduced system. They do not prove that any of
the seven vertices is realizable, transport either invariant through the
Laurent normalization, close $(72,108)$, or raise the degree floor. In
particular, the approximate-root paper proves only an inequality for the
minor-root expression where equality would be needed for the proposed family
exclusion; that route is not used here.

\emph{The transport typing gate (EXP-097).} These original-pair invariants
cannot be imposed directly on the $51$ final Laurent coefficients without
additional boundary data. If $R(x)=\operatorname{Res}_y(F,G)$ and
\(\phi_c(x,y)=(x^{-1},x^cy)\), then [D; exact symbolic controls]
\[
 \operatorname{Res}_y(\phi_cF,\phi_cG)=x^{cpq}R(x^{-1}),
\]
where $p=\deg_yF$ and $q=\deg_yG$. For the final polygons
$(p,q,c)=(16,24,4)$, so $cpq=1536$: an input exponent interval
$[\nu,21]$ becomes $[1515,1536-\nu]$, of width $21-\nu$.
The reduction first swaps the selected coordinate eliminant and then
localizes $x$, making $x^\nu$ a unit; Proposition~4.3 does not retain the
missing $x=0$ divisor order $\nu$. Thus absolute degree $21$ remains a
reconstruction gate for an original polynomial pair, not a direct equation
on the reduced coefficients. Laurent exponent width is a conditional
replacement only after proving $\nu=0$ and swap compatibility. A complete
boundary-divisor ledger could supply that missing theorem; none is claimed
here.

\subsection*{Constructible certificate strata and the first exact chart transition}

The earlier phrase ``finite chart cover by localized certificates'' needs a
commutative-algebra qualification (EXP-098). Suppose principal opens \(D(s_i)\)
cover the parameter space and each carries a localized lift \(c_i\) with
\[
 c_i^TM=0,\qquad c_i^Tb=1.
\]
Clearing denominators gives global syzygies pairing to powers \(s_i^{N_i}\).
Those powers generate the unit ideal, so a polynomial combination is one
global covector pairing to \(1\). A pure principal-open cover of lifted
syzygies is therefore not stronger than the global-covector target.

The strict generalization is a \emph{constructible rank stratification}: after
certifying a generic open, quotient by the residual closed-locus ideal and
recompute the kernel. New syzygies can appear after this non-flat
specialization and need not lift. EXP-098 verifies the distinction on an exact
universally inconsistent control whose global pairing ideal is the proper
ideal \((x)\), while a new certificate appears on \(V(x)\).

This mechanism now reaches the actual GGHV matrix (EXP-099--101). Common-flag
shortcuts fail, but they isolate the first genuine interior interaction:
\[
 s=\varepsilon_{(0,1)},\qquad t=\varepsilon_{(1,7)}.
\]
For the selected \(125\times125\) augmented minor, a Sylvester reduction gives
the exact factorization [MV]
\[
 \frac{\det A(s,t)}{\det A_0}
 =
 \frac{(st-8)^6\bigl(2^{15}s^9-(st-8)^7\bigr)}{2^{39}}.
\]
The mixed coefficient also follows independently from
\[
 \operatorname{tr}B_s=\operatorname{tr}B_t=0,\qquad
 \operatorname{tr}(B_sB_t)=\frac{13}{8}.
\]
At \((s,t)=(-8,-1)\), the first minor vanishes, but exact ranks remain
\[
 \operatorname{rank}M=124,\qquad
 \operatorname{rank}[M\mid b]=125.
\]
Pivot extraction gives an alternative right-hand-side-containing minor
nonzero there. The first two minors cover the component \(st=8\) and leave
only
\[
 2^{15}s^9=(st-8)^7,
\]
with normalization
\[
 s=8u^7,\qquad t=\frac{8u^9+1}{u^7},\qquad u\ne0.
\]
At \(u=1\), both old minors vanish but a third exact augmented minor is
nonzero (EXP-102). Its dense pullback has combined direction rank \(121\) and
hit the declared five-minute budget. EXP-103 replaces that backend by
root-of-unity evaluation and NTT interpolation of complete maximal-minor
polynomials. An endpoint gate catches an 81-degree cancellation that a
modular gcd alone would miss.

EXP-104 resolves that gate over characteristic zero [MV]. The third determinant
is evaluated exactly at 100 integer values. Its assignment bounds are
\([1547,1646]\); after division by \(u^{1547}\), exact degree-99 interpolation
shows that coefficients \(0,\ldots,80\) vanish and coefficient \(81\) does
not. Thus its exact support is \([1628,1646]\). A second row chart has exact
support \([777,903]\).

The support graphs of both charts admit connected row/column weights modulo
9 (EXP-105). With \(z=u^9\), removal of the proved Laurent monomials and
integer contents gives the primitive determinants
\[
 F(z)=21-96z-1024z^2,\qquad G(z)=(8z+1)^{14}.
\]
The only zero of \(G\) is \(z=-1/8\), where \(F(-1/8)=17\). More strongly,
the artifact contains integer polynomials \(A,B\) satisfying
\[
 A(z)F(z)+B(z)G(z)=17^{14}.
\]
Therefore the two normalized maximal minors generate the unit ideal over
\(\mathbb{Q}[u,u^{-1}]\). Together with the first two charts, this proves
rank \(125\) for the augmented matrix at every point of the declared
\(\{(0,1),(1,7)\}\) coefficient slice [MV].

The grading is not confined to that slice. For every nonconstant lower
monomial \(x^py^q\) tested in the persisted 26-parameter family, EXP-106
finds the exact variable-weight law
\[
 w_{p,q}\equiv q-p+1\pmod9.
\]
All 23 remaining nonconstant directions preserve both connected chart
gradings; the constant direction is bracket-zero. This supplies a scalable
next coordinate system, not a full-family proof. The promoted direction
\((0,7)\) has selected ranks 53 and 41; after setting
\(z=u^9\) and \(y=v/u\), both determinant charts have \(z\)-width 14.

EXP-107 reconstructs both bivariate minors on independent 16-by-64 NTT grids
and verifies them by off-grid determinant evaluations [MV]. The endpoint-safe
minor is independent of \(y\), so it remains \(G(z)=(8z+1)^{14}\). The other
minor restricts at \(z=-1/8\) to a squarefree degree-12 polynomial \(Q(y)\).
Thus the first two charts leave only twelve geometric points in that fiber,
not a curve.

EXP-108 selects a third 125-row chart at the residual fiber [MV]. Removing
the connected mod-9 row/column weights turns both fiber matrices into exact
rational matrices. Interpolation under structural degree bounds 14 and 13,
using 29 exact determinants and two independent controls, gives primitive
polynomials \(Q,H\in\mathbb Z[y]\). The polynomial \(Q\) is irreducible over
\(\mathbb Q\), while, up to sign,
\[
\begin{split}
H(y)={}&y(49y^2+168y+192)(2401y^4+1568y^2+1024)\\
&\quad\cdot(117649y^6-403368y^5+921984y^4
-1053696y^3+786432).
\end{split}
\]
Their exact gcd is one. The persisted integer B\'ezout identity has nonzero
constant
\[
3036277895244878564727703434378848855121939007419328699039744.
\]
Since every common zero with \(G\) must lie over \(z=-1/8\), the three minors
prove rank 125 throughout the declared
\(\{(0,1),(1,7),(0,7)\}\) coefficient slice.

These results establish the first complete constructible chart cover on an
actual three-coefficient GGHV slice and expose a global lower-family grading.
They do not cover simultaneous variation of the other coefficients, exclude
the full \((72,108)\) family, or raise the degree floor.

\subsection*{Full-family rank audit and the exact 36-column core}

The next full-family audit corrects a tempting but vacuous rank target. The
\(Q\)-polygon contains the constant monomial, whose column is identically zero:
\[
[P,1]=0.
\]
Hence \(\operatorname{rank}M(\varepsilon)\le124\) for every parameter value.
The exact nonzero pinned augmented minor from EXP-059 then proves that
\(\operatorname{rank}M=124\) and
\(\operatorname{rank}[M\mid b]=125\) generically [D]. This is inconsistency on
a nonempty Zariski-open subset, not on every fiber. The unresolved object is the
common zero locus of augmented maximal minors. EXP-111 also corrects the row
pool: the forced-only construction has 289 rows, while the complete union
inside the same reduced pool has 302 rows [MV].

EXP-112 removes the structural constant column, retains the target column, and
normalizes all 51 exact perturbation matrices on a pinned 125-row chart [MV].
The union dependency graph has one 36-vertex cyclic component, three cyclic
singletons, and 86 acyclic singletons. Consequently the selected determinant
factors as
\[
\det A_{\mathrm{sel}}(\varepsilon)
=\det(A_0)(1+\varepsilon_{(1,0)})^3
\det C_{36}(\varepsilon),
\]
where \(C_{36}\) depends on only 24 parameters. The other 27 directions are
determinant-inert on this chart. On the forced axis,
\[
\chi_{N_{(1,0)}|C_{36}}(\lambda)
=\lambda^{23}(\lambda-1)^{13},
\]
recovering the total factor
\((1+\varepsilon_{(1,0)})^{16}\).

Graph support is not determinant support. EXP-113 finds two
deletion-minimal triples that each make the 36-core strongly connected:
\[
T_A=\{(0,1),(0,7),(2,9)\},\qquad
T_B=\{(0,1),(0,5),(1,0)\}.
\]
EXP-114 computes their exact determinants [MV]. On \(T_A\), the graph-active
\((2,9)\) direction cancels identically, and the degree-21 determinant has only
24 monomials and factors in degrees \(3,12,6\).

On \(T_B\), write
\[
a=\varepsilon_{(0,1)},\qquad
b=\varepsilon_{(0,5)},\qquad
d=1+\varepsilon_{(1,0)}.
\]
Then
\[
\det C_{36}|_{T_B}=2^{-42}G_{54}(a,b,d)H_{63}(a,b,d),
\]
where
\[
\begin{aligned}
G_{54}={}&30720000a^6b^4+48828125a^3b^{11}
+150000000a^3b^8d\\
&+64000000a^3b^5d^2-39321600a^3b^2d^3
+16777216d^6,
\end{aligned}
\]
and
\[
\begin{aligned}
H_{63}={}&4096a^9+184320a^6bd^2-1800000a^3b^5d^3
+1843200a^3b^2d^4\\
&+1953125b^9d^4+3000000b^6d^5+1536000b^3d^6
+262144d^7.
\end{aligned}
\]
These factors are weighted homogeneous for
\(\operatorname{wt}(a,b,d)=(7,3,9)\), of weighted degrees 54 and 63.
They expose a weighted open/boundary split for the next chart calculation:
normalize \(d\ne0\), then treat \(d=0\) with alternative minors from the
complete 302-row system. No result in this subsection excludes the residual
factor loci, the full \((72,108)\) family, or \(JC(2)\).

\subsection*{Weighted residual transitions and the \(d=0\) quotient}

EXP-115 first normalizes the open chart \(d\ne0\) by
\[
a=A u^7,\qquad b=B u^3,\qquad d=u^9.
\]
Writing \(X=A^3\), \(G_{54}(A,B,1)\) remains irreducible, while
\[
H_{63}(A,B,1)=L(A,B)Q(A,B)
\]
splits into irreducible degree-three and degree-six factors over
\(\mathbb Q[A,B]\) [MV]. Exact good-prime alternative minors prove that none
of \(G_{54},L,Q\) is wholly contained in the rank-deficient locus. The
\(L\)-transition also lifts to the rational point
\((A,B,d)=(0,-4/5,1)\). Proper intersections on these components remain.

The boundary \(d=0\) is structural. Cancelling the forced \(x\)-term puts
\[
P=a y+b y^5+y^8(1-xy)^8
\]
inside the retained \(Q\)-space, so \([P,P]=0\) is an exact polynomial right
kernel [D]. Its \(y^8\)-coordinate is the unit \(1\); deleting that column
gives a 302-by-124 quotient augmented matrix with 123 coefficient columns.
The origin has exact rank profile \(112/113\), while the nonzero controls
\((0,1)\), \((1,0)\), and \((-9,12/5)\) have profile \(123/124\) [MV].

EXP-116 compresses the quotient graph into cyclic blocks of sizes
\[
51,11,10,9,8,7
\]
and sixteen cyclic singletons. EXP-117 computes the separately budgeted
51-block and all smaller determinants. In original coordinates the selected
quotient determinant is, up to a persisted nonzero rational scalar,
\[
\det A_{\mathrm{quot}}
\doteq
b^{32}F_{28}(a,b)
\prod_{c\in\mathcal C}(ca^3+78125b^7),
\qquad
\mathcal C=\{49152,32768,16384,8192,4096\},
\]
where
\[
\begin{aligned}
F_{28}={}&13750615562268966912a^{12}
-389229646907965440000000a^9b^7\\
&+6149594411200000000000000a^6b^{14}\\
&+865348510742187500000000a^3b^{21}\\
&-71190297603607177734375b^{28}.
\end{aligned}
\]
All factors are homogeneous for
\(\operatorname{wt}(a,b)=(7,3)\). On \(b\ne0\), set \(z=a^3/b^7\). The
residual becomes one squarefree degree-nine equation: five rational roots
\[
-\frac{78125}{4096},\quad -\frac{78125}{8192},\quad
-\frac{78125}{16384},\quad -\frac{78125}{32768},\quad
-\frac{78125}{49152},
\]
and four roots of the quartic obtained from \(F_{28}/b^{28}\) [MV].
Thus the selected quotient residual is reduced to \(b=0\) and nine invariant
points.

EXP-118 supplies the exact alternative chart [MV]. Its complete quotient
determinant is
\[
\det A_{\mathrm{alt}}\doteq a^{107}.
\]
The persisted scalar is a nonzero integer. This chart covers \(a\ne0\),
including all nine invariant points. On \(a=0\), the selected determinant is
exactly \(b^{95}\), and at the origin the quotient rank profile is
\(112/113\). These three cases close the complete \(d=0\) quotient plane.
This is a boundary-stratum closure, not a result about \(d\ne0\).

\subsection*{Exact weighted-open charts and component ideals}

EXP-119 verifies the full 302-by-125 augmented matrix is covariant for
\(\operatorname{wt}(a,b,d)=(7,3,9)\) [MV]. On \(d=1\), a first alternative
row basis gives an exact determinant with 114 monomials. Writing \(X=A^3\),
it factors as a nonzero rational scalar times
\[
B^{36}K(X,B)
\prod_{c\in\{4096,8192,16384,32768,49152,65536,86016\}}
(cX+78125B^7),
\]
where \(K\) has degree 6 in \(X\), degree 22 in \(B\), and 30 monomials.
Its gcd with each irreducible selected component \(G,L,Q\) is one.
Nevertheless, the three exact elimination resultants are nonconstant, of
degrees 852, 324, and 648 in \(B\). Hence this chart removes every curve
component but leaves three finite proper intersections.

EXP-120 evaluates the independently selected \(G\)-basis [MV]. The first
rational normalization \((A,B,d)=(1,0,1)\) works, the largest cyclic block
has size 25, and the exact determinant has 21 monomials. Up to a nonzero
rational scalar its invariant form is
\[
\Delta_G=X^{30}L(X,B)Q(X,B)R(X,B),
\]
where
\[
R=-156250B^7+15625B^6X+28000B^3X+3200B^2X^2+1024X.
\]
Factorwise exact Groebner calculations give
\[
(G,\Delta_{LQ},L)=(G,\Delta_{LQ},X)
=(G,\Delta_{LQ},Q)=(G,\Delta_{LQ},R)=(1).
\]
Therefore the three charts cover the complete \(G\) component [D].

The corresponding EXP-120 ideals on \(L\) and \(Q\) are nonunit but
zero-dimensional. The \(L\) lexicographic eliminant has degree 108 in \(B\)
and squarefree degree 73. For \(Q\), a six-element graded-lex basis proves
zero-dimensionality; the optional FGLM conversion reached its declared
180-second gate, so no lex eliminant or point count is claimed. The explicit
\(LQ\) factor in \(\Delta_G\) shows that this basis cannot remove either
remaining component. The next exact charts must be selected directly on
those finite residual schemes.

EXP-121 performs that selection [MV]. At \(p=1013\), thirteen affine points
of the \(L\) residual have coefficient/augmented rank profile \(124/125\);
at \(p=1033\), eighteen points of the \(Q\) residual have the same profile.
Deterministic row extraction on the two components returns the same
125-row basis, although it differs from the pinned basis by 68 rows. Its
rational anchor is \((A,B,d)=(1,0,1)\), its largest cyclic block has size
26, and five direct 125-by-125 controls agree with the block product. The
exact determinant has total degree 108 and 23 monomials:
\[
 \Delta_{\mathrm{fin}}(A,B)=A^{87}R_{\mathrm{fin}}(A^3,B),
 \qquad
 \deg_XR_{\mathrm{fin}}=5,\quad
 \deg_BR_{\mathrm{fin}}=18.
\]

For \(L\), the exact graded-reverse-lexicographic calculation returns
\[
 (L,\Delta_{LQ},\Delta_{\mathrm{fin}})=(1).
\]
For \(Q\), a raw three-generator calculation exceeded its declared gate, so
the persisted certificate uses the exact factor
\(\Delta_{LQ}=B^{36}\Delta^\circ\). On \(B=0\),
\[
 \gcd\!\left(256X^2-1024X+4096,\ X^{32}\right)=1.
\]
On \(\Delta^\circ=0\), reduction modulo the quadratic \(Q\) makes both
\(\Delta^\circ\) and \(\Delta_{\mathrm{fin}}\) linear in \(X\). The
exceptional coefficient gcd is one, and the resulting degree-144 and
degree-176 compatibility polynomials in \(\Q[B]\) also have gcd one.
Consequently
\[
 (Q,\Delta_{LQ},\Delta_{\mathrm{fin}})=(1).
\]
This is an exact zero-set split and univariate gcd certificate, not a
modular extrapolation.

Together, EXP-118--121 close the complete three-parameter \(T_B\)
restriction: EXP-118 covers \(d=0\), EXP-120 covers \(G\) on \(d\ne0\), and
EXP-121 covers the remaining finite \(L/Q\) strata. They do not close the
24-parameter core, the full 51-parameter family, \((72,108)\), the planar
degree floor, or \(JC(2)\).

\subsection*{First higher-dimensional lift from the cyclic core}

EXP-122 restores the other 21 directions in the 24-parameter cyclic core on
the shared EXP-121 basis [MV]. No restored direction is
determinant-inert on its complete anchor line. Thirteen directions act inside
the existing size-26 cyclic block, 16 enlarge the fixed-\(d\) SCC, and every
restored direction contributes at first order or pairwise mixed order. The
free-lift prediction is therefore refuted for this basis.

The eight directions \((2,j)\), \(2\leq j\leq9\), nevertheless form a sparse
class: each enlarges the SCC but has a linear anchor-line determinant.
Direction \((2,9)\) has the smallest selected rank, normalized support, and
enlarged SCC. EXP-123 computes its complete symbolic lift on \(d=1\) [MV].
Writing \(C=\varepsilon_{(2,9)}\),
\[
 X=A^3,\qquad Y=A^2C,
\]
the shared determinant is
\[
 \Delta_{(2,9)}(A,B,C)
 =A^{87}\bigl(R(X,B)+Y S(X,B)\bigr).
\]
Here \(R\) is the 23-term EXP-121 polynomial, \(S\) has 18 terms, and
\[
 \gcd_{\Q[X,B]}(R,S)=1.
\]
Two modular probes precede the exact computation. The claim rests on the
characteristic-zero determinant, its \(C=0\) regression, and four direct
full-matrix rational controls.

Consequently, on \(A\ne0\) and \(S\ne0\), the selected exceptional locus is
the rational graph
\[
 Y=-\frac{R(X,B)}{S(X,B)}.
\]
The remaining base locus \(V(R,S)\subset\mathbb A^2_{X,B}\) is
zero-dimensional because \(R\) and \(S\) are coprime. This is the first
higher-dimensional constructible reduction from the 24-parameter core [D].
It is not a four-parameter chart cover: an alternative minor must still cover
the rational graph, the finite base locus, and the separate divisor \(A=0\).
The complete \((72,108)\) case, the planar degree floor, and \(JC(2)\) remain
open.

\subsection*{Recursive atlas on the rational graph}

EXP-124 selects an alternative complete-row basis directly on the graph
[MV]. The same one-row replacement is full rank over two good primes, and
its exact determinant is
\[
 \Delta_{\rm alt}(A,B,C)=A^{90}N(A^3,B).
\]
It is independent of \(C\). The 21-term polynomial \(N(X,B)\), of total
degree 16, factors over \(\Q\) as a nonzero scalar times
\[
 N(X,B)=F_3(X,B)F_6(X,B)F_7(X,B),
\]
where the factors have total degrees \(3,6,7\) and are each coprime to
\(R,S\). Thus this minor covers the dense graph open \(N\ne0\), reducing the
positive-dimensional residual to three explicit plane curves.

EXP-125 recurses first on
\[
 F_3(X,B)=(5B+4)^3+16X.
\]
Prime availability is not automatic: the originally named primes were
refuted by cubic-residue and factor-point obstructions, and an intermediate
audit that allowed arbitrary \(X\) was also rejected because \(X=A^3\) is
required. The final cube-locus audit selected 739 and 811. On both primes,
four points on each of \(F_3,F_6,F_7\) have coefficient/augmented rank
\(124/125\), and each curve supplies a new cross-prime basis [MV].

The exact \(F_3\)-selected determinant has one cyclic block of size 32. After
restriction to \(Y=-R/S\) and substitution
\(X=-(5B+4)^3/16\), its univariate numerator factors as
\[
 U_3(B)\doteq(5B+4)Q_6(B)^2Q_9(B)Q_{15}(B).
\]
Here \(\doteq\) denotes equality up to a nonzero rational scalar and the
subscripts are exact irreducible degrees. The linear factor forces \(A=0\);
\(Q_6\) divides both \(R|_{F_3}\) and \(S|_{F_3}\); and \(Q_9,Q_{15}\)
divide neither. Consequently, on \(AS\ne0\), only
\[
 Q_9(B)Q_{15}(B)=0
\]
remains on \(F_3\): 24 normalized parameter values and 72 lifted points over
an algebraic closure [MV/D]. This is a finite reduction of one component,
not a complete four-parameter cover.

EXP-126 treats a maximal minor as a section on the next residual curve
[MV]. The persisted \(F_6\) basis is exactly the \(F_3\)-selected basis
whose characteristic-zero determinant was already reconstructed. After
hash, anchor, SCC, graph-numerator, modular-sample, and direct-determinant
checks, the graph numerator has the nonzero class
\[
 h_6(X,B)=U(B)X+V(B)
 \quad\text{in}\quad \Q[X,B]/(F_6).
\]
The function-field norm is computed independently by a Sylvester resultant
and by the determinant of multiplication by \(h_6\) in the quadratic
quotient. It has degree 74 and factors, up to scalar, as
\[
 \operatorname{Norm}(h_6)
 \doteq D_2D_3^4D_6^2Q_{18}Q_{30}.
\]
The five factors are irreducible, with subscripts denoting degrees. Direct
arithmetic in the corresponding quotient fields shows that \(D_2\) gives
\(X=S=0\), while \(D_3,D_6\) give \(R=S=0\). The degree-3 case is a double
\(F_6\) fibre on which \(U=V=0\), so it is classified at the double root
rather than by a projection resultant alone. Neither \(Q_{18}\) nor
\(Q_{30}\) meets \(X=0\) or \(S=0\). Consequently, on \(AS\ne0\), the
effective residual is
\[
 Q_{18}(B)Q_{30}(B)=0.
\]
This gives 48 normalized values and 144 lifted algebraic points [MV/D].
It is the first explicit use of the maximal-minor / curve-section /
zero-divisor / function-field-norm dictionary in this atlas.

The graph curve \(F_7\), both finite \(F_3/F_6\) residuals, the
finite base locus, and the separate divisor \(A=0\) remain, as do the full
\((72,108)\) case, the planar degree floor, and \(JC(2)\).

\emph{The frontier, corrected.} Beyond $125$, our first probe case (the chain
$(8,40)/(8,28)/(11/4,7)$, max degree $144$) turns out already excluded in the
published literature: the closing remark of \cite{GGV2017} states that the corner
data $(8,28),(8,40)$ forces $A_0' = (8,4)$, which is not a possible last lower
corner. We had not fully mined that remark; several of the twenty-four
configurations in $[125,150]$ may be similarly settled, and a re-audit against the
excluded family $\wp(n', n'-1)$ precedes any new attack (EXP-082, EXP-083).

\subsection*{The counterexample as an incidence map, and its exclusion in the plane}

A complementary lens sharpens why the two-variable case resists the very mechanism
that settled dimension three. The 2026 counterexample is not an arbitrary polynomial
map: it is the coefficient map of the product of a linear and a quadratic binary
form, restricted to a normalized slice. Concretely, on
$V_1\times V_2\to V_3\times\mathbb{A}^1$, $(L,Q)\mapsto(LQ,\operatorname{Res}(L,Q))$,
the slice $X_H=\{(L,Q):\operatorname{Res}(L,Q)=1,\ [LQ]_{\text{next}}=1\}$ carries the
projection $(L,Q)\mapsto LQ$, which is étale on the coprime (nonzero-resultant) locus
by a one-line argument: multiplication forgets exactly the relative scaling
$(L,Q)\mapsto(\lambda L,\lambda^{-1}Q)$, and the resultant, of weight $\deg Q-\deg L$
under that scaling, detects it. The projection is generically $n$-to-one (the choices
of linear factor), hence not injective. We verified [MV] the full identification: the
displayed counterexample equals this coefficient map up to linear normalization, with
$X_H\cong\mathbb{A}^3$. So, for this family, the entire Jacobian question reduces to a
single \emph{recognition of affine space}: is $X_H\cong\mathbb{A}^n$?

In dimension three the recognition holds by a toric coincidence: the residual
$\mathbb{G}_m$-grading has weights $(1,-1,-2)$, and the associated semigroup generators
form a unimodular lattice basis, so $X_3\cong\mathbb{A}^3$. In the plane it \emph{fails},
and we can say why in one invariant. The controlling quantity is
$n-2=\deg Q-\deg L$, the resultant weight under the relative scaling; it is
simultaneously the degree of the Danielewski core
$\{x^{\,n-2}W=B^{\,n-2}-1\}$ and the number of boundary components. For $n=2$ it is
$0$: the resultant becomes scaling-invariant, unable to rigidify the normalization,
and the slice degenerates (reducible; equivalently, the finite-root model is
$\mathbb{A}^2$ minus the discriminant, which carries a nonconstant unit and is
therefore not $\mathbb{A}^2$). For $n=3$ it is $1$, the unique value that is both
nonzero and linear, giving $X_n\cong\mathbb{A}^n$; for $n\ge4$ it is $\ge2$, and the
class group obstructs. Thus the exact structure of the known counterexample cannot
occur in the plane [MV], a conclusion consistent with the companion's
$\mathbb{G}_m$-equivariant rigidity theorem, which forbids the torus-graded planar
analog independently. This excludes the incidence/root-selection class, the structure
of every constant-Jacobian noninvertible map known to date; it is strong structural
evidence for $JC(2)$, not a proof, since a planar counterexample of some entirely
different structure is not addressed.

\section{Positioning and novelty}\label{sec:positioning}

Per the adversarial novelty pass (same dossier), no prior statement of
Theorems~\ref{thm:t1} or \ref{thm:t2} was found. The novelty record for the companion's
equivariant rigidity theorem is updated: Shaska \cite{Shaska2026}, submitted on 2026-07-22
after EXP-010 was declared, independently proves the full all-signature classification.
Same-sign actions may give nonlinear triangular automorphisms; opposite-sign actions give
the linear case of EXP-010. No priority or novelty claim is made for that classification.
Theorem~\ref{thm:t3} was not found as stated, but its
operator identity is within the standard bracket calculus \cite{Joseph1975, GGV2017}: the
statement appears new, the method is classical. Theorem~\ref{thm:t4}'s counterexample half
is essentially contained in \cite{GGV2017} (Prop.~4.1) with van den Essen's subrectangular
normalization \cite{vdEssen2000}; we claim only the sharp dichotomy form. All proofs are
elementary once seen; a folklore risk remains and is recorded, with the primary-source
audit trail (including its open verification items) persisted in the repository.

\section{Program}

Priorities, from the standing routes evaluation: (1) promote the
recursive graph atlas from the \(F_3/F_6\) finite reductions to an exact
divisor/norm cover of \(F_7\), followed by the finite graph/base-locus points
and the separate \(A=0\) boundary; (2) retain the complete EXP-118--121 \(T_B\) cover as a
regression gate rather than adding redundant slice charts; (3) retain the
exact EXP-108 lift through \((0,6)\) as a bounded chart control rather than
a terminating programme; (4) reopen the intersection-$21$ transport only
after constructing a complete boundary-divisor ledger through the swap,
Laurent cuts, and final inversion; (5) pursue Newton resolution only after deriving a
condition beyond the already-retained first $D=72$ branch; and (6) retain the
properness reformulation ($JC(2)$ holds iff every planar Keller map has empty
Jelonek set; exact certificates shipped, EXP-014) as a cross-check instrument.
The reproducibility contract: every [MV] claim has a persisted script and
output; hypotheses are declared before runs; refuted predictions remain in the
record.

\small
\begin{thebibliography}{99}
\bibitem{Alpoge2026} L. Alp\"oge, announcement of the counterexample, X (@\_\_alpoge\_\_),
2026-07-19/20; map produced with Claude Fable 5 (Anthropic); question posed by Akhil.
\bibitem{Magnus1955} A. Magnus, On polynomial solutions of a differential equation,
\emph{Math. Scand.} \textbf{3} (1955), 255--260. (The gcd $=1$ theorem; Magnus' formula is
in \emph{Proc. Amer. Math. Soc.} \textbf{5} (1954), 256--266.)
\bibitem{AppelgateOnishi1985} H. Appelgate, H. Onishi, The Jacobian conjecture in two
variables, \emph{J. Pure Appl. Algebra} \textbf{37} (1985), 215--227.
\bibitem{Nagata1989} M. Nagata, Some remarks on the two-dimensional Jacobian conjecture,
\emph{Chinese J. Math.} \textbf{17} (1989), 1--7.
\bibitem{Moh1983} T.-T. Moh, On the Jacobian conjecture and the configurations of roots,
\emph{J. Reine Angew. Math.} \textbf{340} (1983), 140--212.
\bibitem{Heitmann1990} R. Heitmann, On the Jacobian conjecture, \emph{J. Pure Appl.
Algebra} \textbf{64} (1990), 35--72.
\bibitem{Abhyankar2008} S. S. Abhyankar, Some thoughts on the Jacobian conjecture, Parts
II--III, \emph{J. Algebra} \textbf{319} (2008), 1154--1248; \textbf{320} (2008),
2720--2826.
\bibitem{GGV2017} J. A. Guccione, J. J. Guccione, C. Valqui, On the shape of possible
counterexamples to the Jacobian conjecture, \emph{J. Algebra} \textbf{471} (2017), 13--74;
arXiv:1401.1784.
\bibitem{GGHV2017} J. A. Guccione, J. J. Guccione, R. Horruitiner, C. Valqui, Some
algorithms related to the Jacobian Conjecture, arXiv:1708.07936.
\bibitem{GGHV2022} J. A. Guccione, J. J. Guccione, R. Horruitiner, C. Valqui, Increasing
the degree of a possible counterexample to the Jacobian Conjecture from 100 to 108,
arXiv:2204.14178.
\bibitem{GGHVApproxRoots} J. A. Guccione, J. J. Guccione, R. Horruitiner, C. Valqui,
Approximate roots and intersection numbers, arXiv:1708.09367v2.
\bibitem{LeeLi2024} K. Lee, L. Li, On the two-dimensional Jacobian conjecture:
Magnus' formula revisited IV, arXiv:2408.01279.
\bibitem{MakarLimanovTrakhtenberg2024} L. Makar-Limanov, L. Trakhtenberg,
Properties of a Jacobian mate, MPIM Preprint 2024 (33);
\href{https://doi.org/10.1007/s40863-025-00520-4}{doi:10.1007/s40863-025-00520-4}.
\bibitem{Joseph1975} A. Joseph, The Weyl algebra: semisimple and nilpotent elements,
\emph{Amer. J. Math.} \textbf{97} (1975), 597--615.
\bibitem{vdEssen2000} A. van den Essen, \emph{Polynomial Automorphisms and the Jacobian
Conjecture}, Progress in Mathematics 190, Birkh\"auser, 2000.
\bibitem{Miyanishi2023} M. Miyanishi, Equivariant Jacobian Conjecture in dimension two,
\emph{Transform. Groups} \textbf{28} (2023), 951--971.
\bibitem{DuboulozPalka2018} A. Dubouloz, K. Palka, The Jacobian Conjecture fails for
pseudo-planes, \emph{Adv. Math.} \textbf{339} (2018), 248--284.
\bibitem{Shaska2026} T. Shaska, Graded Keller maps and the Jacobian Conjecture,
arXiv:2607.20210v1, 2026.
\end{thebibliography}
\normalsize

\end{document}
